\section{Finite certificates and reproducibility}\label{sec:certificates}

The electronic appendix contains the fixed rational inputs, verification
programs, a file index, and a manifest of SHA-256 hashes. The programs
recompute the inequalities from these inputs by the methods below.

\subsection{Elementary enclosures}

All rational calculations use integer numerators and positive integer
denominators. For a positive integer $m$, a square-root enclosure is
obtained by taking the integer square root of $\lfloor m^2a/b\rfloor$
for a rational $a/b\ge0$;
division by $m$ gives the lower bound and addition of $1/m$ gives
an upper bound unless the square is exact. Thus the directions of
the rounded radical bounds are determined by integer inequalities.

The enclosures $6283/2000<333/106<\pi<22/7$ follow from Machin's identity
\[
 \pi=16\arctan(1/5)-4\arctan(1/239).
\]
Indeed $\tan(4\arctan(1/5))=120/119$, and subtraction of
$\arctan(1/239)$ gives tangent one with angle in $(0,\pi/2)$.
If $A_m(x)=\sum_{k=0}^{m-1}(-1)^kx^{2k+1}/(2k+1)$, the
alternating-series bounds give
\[
 16A_6(1/5)-4A_3(1/239)<\pi
 <16A_7(1/5)-4A_2(1/239).
\]
Both comparisons with the displayed rational endpoints are exact.
Finer enclosures are obtained by retaining more terms.

For $x\ge1$, first divide $x$ by a power of two to put it in $[1,2)$.
With $t=(x-1)/(x+1)$ on this interval,
\[
 2\sum_{j=0}^{m-1}\frac{t^{2j+1}}{2j+1}
 \le\log x
 \le2\sum_{j=0}^{m-1}\frac{t^{2j+1}}{2j+1}
       +\frac{2t^{2m+1}}{(2m+1)(1-t^2)}.
\]
The same formula at $x=2$ encloses the range-reduction contribution.
The formulas for $\log(1\pm r)$ in
\eqref{mf:directed-log} retain their complete alternating or
positive geometric remainder. For an integer $m\ge0$, put
$E_m(x)=\sum_{j=0}^m x^j/j!$. If $0\le x<m+2$, then
\begin{equation}\label{ce:exp}
 E_m(x)\le e^x\le E_m(x)
                +\frac{x^{m+1}}{(m+1)!(1-x/(m+2))}.
\end{equation}
Negative exponentials are bounded by reciprocals of positive
enclosures. The seven rows in Proposition~\ref{ar:energy}, for
example, use the lower logarithm bound with $48$ terms and
$e^{-x}\le1/E_{160}(x)$. Their checker uses only rational arithmetic.

The interval computations use directed real and complex balls.
An asserted strict inequality means that the entire resulting
enclosure has the required sign. Input floating-point test points are converted to
their exact binary rationals before use.

\subsection{Entire integrals}

The only numerical integrations in the field certificates are of
entire functions of the form $\exp(P(t))$, with $P$ a polynomial,
over a finite real interval. The preceding analytic lemmas bound the
finite-product errors separately from these model integrals.

There is a direct series verification of each such integral. After
an affine change of variable, write
$P(t)=p_0+\sum_{j=1}^d p_jt^j$ on $[0,1]$ and
$\exp(P(t)-p_0)=\sum_{k\ge0}a_kt^k$.
The coefficients are determined by
\[
 a_0=1,\qquad
 ka_k=\sum_{j=1}^{\min(d,k)}jp_ja_{k-j}.
\]
For any rational $R>1$, Cauchy's estimate gives
$|a_k|\le M_RR^{-k}$, where
$M_R=\exp(\sum_{j=1}^d|p_j|R^j)$.
Consequently
\[
 \left|\int_0^1e^{P(t)}\,dt
       -e^{p_0}\sum_{k=0}^{m}\frac{a_k}{k+1}\right|
 \le |e^{p_0}|
       \frac{M_RR^{-(m+1)}}{(m+2)(1-1/R)}.
\]
Subdividing the interval, or increasing $m$, gives arbitrarily
small error bounds. Formula \eqref{ce:exp} encloses $M_R$.
The field checkers also supply validated analytic quadrature
using Arb/Acb; all integrands satisfy its entire-function
hypothesis. The separate series implementations in the
electronic appendix give the coefficient recurrence and the
Cauchy remainder explicitly.

For the Gaussian shells in Proposition~\ref{mf:affine}, the rational
integration-by-parts and Mills bounds in its proof suffice.
An additional verifier evaluates the Gaussian tails by the
complementary error function with interval bounds.

\subsection{Polynomial bounds on whole boxes}

The companion centers are defined by finite recurrences, so their
coefficient identities are checked by exact polynomial arithmetic.
For the Gaussian center, the moments
at $s=\pm2\pi i$ obey $I_0=0$ and
$I_m=(1-mI_{m-1})/s$. This derives the center and residual
coefficients exactly. The identity $e^s=1$ is used only at
$s=\pm2\pi i$; the resulting polynomial bounds may then be
verified on the larger real interval $6\le |s|\le44/7$.

After rescaling each real variable to $[0,1]$, a polynomial of
degree at most $d$ has the Bernstein representation
\[
 \sum_{j=0}^dc_jt^j
 =\sum_{k=0}^d\left(\sum_{j=0}^k
           \frac{\binom{k}{j}}{\binom{d}{j}}c_j\right)
          \binom{d}{k}t^k(1-t)^{d-k}.
\]
The basis functions are nonnegative and sum to one. Therefore
the extrema of the coefficient list bound the polynomial
on the whole interval. Tensor products give the same result
on a box. Exact dyadic subdivision and the de Casteljau
recurrence preserve these bounds. Bounds on complex moduli
are proved by applying this rule to the difference between
a proposed squared bound and the squared modulus.
The common coefficient-domain verifier checks all twenty
Gaussian polynomial bounds by this method.

Where a one-variable polynomial sign is sharper than its
Bernstein enclosure, the programs use its exact Sturm chain.
Polynomial division is over the rationals; signs at both
endpoints are checked separately. The variation difference
counts roots in the open interval. Combined with a sign at
one point and the endpoint checks, this proves the desired
sign on the whole closed interval.

\subsection{Coverage and degree dependence}

The rectangular covers in this paper use either a rational
edge decomposition or a dyadic tree. In the former, checking
every open slab between consecutive rational edges and every
edge itself proves closed coverage; the horizontal intervals
are merged in sorted order. In the latter, every deepest
pair of leaves must be the two halves of the same coordinate
at a common parent. Replacing such a pair by its parent
must eventually produce the complete root box.

All inequalities involving the degree are extended from
$N_0=6250$ by the stated analytic monotonicities. One
frequent step is worth making explicit. If
$H>(N-c)/F$, with $F>0$ and $0\le c<N_0$, then
$H/N>(1-c/N)/F$ increases as a lower bound with $N$.
Together with $Q\ge H^2/N$ this gives, for every $N\ge N_0$,
\[
 Q>\frac{(N_0-c)^2}{N_0F^2}.
\]
This is the energy lower bound used in the product estimates.

The principal finite covers have the following sizes.
The count of contours includes complete boundaries, with
the inverse-map and amplitude comparisons both checked.
\begin{center}\small
\begin{tabular}{lrr}
\toprule
Component&closed cells or rows&additional finite objects\\
\midrule
Affine field&$275+28$&$2750$ paired shells\\
$D\ge25$ transition&$50614$&$186477$ point obstructions\\
$D\ge30$ transition&$26397$&$82823$ point obstructions\\
$D\ge100$ transition&$9745$&$8160$ point obstructions\\
Origin base&$9$ path ranges&$55$ inverse contours\\
Origin $T=1$ strip&$63$&$3831$ inverse contours\\
Origin $X=1$ strip&$301$&$14891$ inverse contours\\
Fine small-witness range&$1$&$41$ inverse contours\\
Low and middle small-witness range&$23$&$1809$ inverse contours\\
High small-witness range&$20$&$411+20$ contour constructions\\
Moment extension&$20$&$954$ inverse contours\\
High-$X$ extension&$10$&$709$ inverse contours\\
Boundary and auxiliary strips&$8$&whole-strip inequalities\\
Compact range&$13$&$2918+2$ inverse contours\\
\bottomrule
\end{tabular}
\end{center}

The file \cert{verification/INDEX.json} maps each mathematical
component to its fixed data and executable checker. The entry
point \cert{verification/run_all.py} runs the included
verifications and writes new results to a separate results
directory. The input hashes are checked before execution.
The programs use Python's exact rational arithmetic, SymPy
for exact polynomial arithmetic, and, where indicated,
\cert{python-flint} for rigorous balls.
The LaTeX source, all fixed inputs, and the verification code are included.
